\section{Introduction}\label{sec:intro}

Consider a network packet-routing function deployed on a server
handling millions of requests per day.  The function computes a hash
of each packet's header and maps it to one of $k$ queues; a collision
sends two packets to the same queue, degrading throughput.  The
production traffic distribution is far from uniform---certain header
patterns are orders of magnitude more common than others.  The
engineering team knows that edge-case collisions exist, because a
fuzzer found one.  What they need to know is the \emph{rate}: what
fraction of production inputs actually trigger the defect?  A single
counterexample does not answer this question.

This scenario is an instance of \emph{distribution-aware reliability
  estimation}.  Given a deterministic program $\prog : \inX \to \inY$,
a discrete distribution $\dist$ over $\inX$, and a Boolean property
$\prop : \inY \to \{0,1\}$, the \emph{operational reliability} of
$\prog$ under $\dist$ with respect to $\prop$ is
\begin{equation}\label{eq:mu-def}
  \mutrue \;=\; \Pr_{x \sim \dist}\!\big[\prop(\prog(x)) = 1\big].
\end{equation}
Reliability in this sense arises in \emph{operational testing}, where
a developer must certify that a production input distribution produces
a failure rate below a contractual threshold; in \emph{probabilistic
  program
  verification}~\cite{geldenhuys2012pse,filieri2015tacas}, where the
question is inherently quantitative; and in \emph{distribution-aware
  property-based
  testing}~\cite{maciver2019hypothesis,claessen2000quickcheck}, where
quantitative coverage supplements counterexample search.

\paragraph{The challenge.}
Three difficulties set reliability estimation apart from conventional
program analysis.

\emph{The answer is a real number, not a Boolean.}  Standard
verification asks $\forall x\colon \prop(\prog(x))$.  Reliability asks
for the \emph{measure} of the failure set $\{x : \prop(\prog(x))=1\}$
under $\dist$.  Counterexample-driven techniques produce a single
witness; bounding the probability of the entire set requires a
different methodology.

\emph{The rare-event regime matters most and is hardest.}  When
$\mutrue \ll 1$, plain Monte Carlo requires $\Theta(1/(\mutrue
\varepsilon^2))$ samples to certify a half-width of $\varepsilon$---a
prohibitive budget when $\mutrue = 10^{-4}$ and $\varepsilon =
10^{-5}$.  Yet rare events are precisely where operational reliability
matters: a defect that occurs once per million inputs under a uniform
distribution may occur once per hundred inputs under production
traffic.

\emph{The certificate is the deliverable.}  A point estimate for
$\mutrue$ without an accompanying confidence interval is operationally
useless: one cannot determine whether $\hat{\mu} = 0.003$ is
consistent with a contractual bound of $\le 0.001$ unless the
estimator's uncertainty is quantified.  Sound, tight intervals---not
just point estimates---are the engineering artifact that reliability
estimation must produce.

\paragraph{Existing approaches.}
Two families of techniques address this problem, each with
complementary strengths and limitations.

\emph{Probabilistic symbolic execution}
(PSE)~\cite{geldenhuys2012pse,filieri2015tacas,borges2014compositional}
enumerates program paths, encodes each path condition as a formula,
and computes its exact probability via model
counting~\cite{chakraborty2013approxmc,yang2023approxmc6}.  PSE is
exact where it applies.  However, path explosion forces incompleteness
on programs with unbounded loops or non-linear arithmetic, and
non-uniform input distributions require ad-hoc extensions of the
model-counting backend that are not always available.

\emph{Plain sampling}~\cite{sampson2014passert,kwiatkowska2011prism}
draws iid inputs, executes the program, and applies a concentration
inequality to the empirical failure rate.  Sampling is robust to all
sources of PSE difficulty, but inherits the $\Theta(1/\varepsilon^2)$
rate of independent estimation.  In the rare-event regime, the sample
budget required to reach a useful half-width is simply too large.

\paragraph{Our approach.}
We propose a middle course between exact symbolic enumeration and
unguided sampling.  \toolname{} maintains a \emph{frontier}: a tree whose leaves
partition the input domain $\inX$ into path-condition regions.  At
each step the algorithm chooses between two actions:
\textsc{Sample}$(\pi)$, which draws additional executions from a
region $\pi$ to tighten its statistical confidence bound, and
\textsc{Refine}$(\pi, b)$, which invokes an SMT solver to split $\pi$
along a branch predicate $b$ harvested from observed executions,
potentially certifying that an entire child region is all-pass or
all-fail under $\prop$.  The choice is governed by a single
gain-per-cost criterion applied uniformly to both action types.

The key insight is a two-part decomposition of the certified
half-width:
\begin{equation}\label{eq:hw-decomp}
  \Pr\!\big[\,|\muhat - \mutrue| \;\le\; \epsstat + \Wopen\,\big]
  \;\ge\; 1 - \delta,
\end{equation}
where $\epsstat$ aggregates per-region statistical uncertainty and
$\Wopen$ is the total probability mass of regions not yet
symbolically certified.  Every \textsc{Sample} action reduces
$\epsstat$; every successful \textsc{Refine} reduces $\Wopen$.  The
algorithm terminates as soon as $\epsstat + \Wopen \le \varepsilon$.
There are no phases, no hand-tuned switching heuristics: the
decomposition~\eqref{eq:hw-decomp} is simultaneously the correctness
argument and the scheduler.

Two design choices distinguish \toolname{} from prior work.  First,
the per-region statistical bound instantiates the predictable-plug-in
empirical-Bernstein (PrPl-EB) confidence sequence of Waudby-Smith and
Ramdas~\cite{waudbysmith2024betting}, which is closed-form,
variance-adaptive, and \emph{anytime-valid}: it remains sound under
\toolname{}'s data-dependent stopping rule, strictly tightening the
union-bounded Wilson construction used in earlier statistical model
checking
tools~\cite{boyer2013plasma,hartmanns2014modest}.  Second, the
closure rule is single-branch: a region is certified if and only if
the SMT backend proves path-determinism ($F_\pi \wedge
\neg\,\text{path}_\pi$ is \unsat{}).  There is no sample-based
fallback.  The price is wider intervals on programs whose closure is
symbolically intractable; the gain is a one-lemma soundness contract
and the elimination of a class of silent bias that affects prior
prototypes.

\paragraph{Contributions.}
This paper makes three contributions.

\begin{enumerate}
\item \emph{The \toolname{} algorithm and its correctness proof.}  We
  introduce \toolname{} and prove that its reported interval $[L,U]$
  covers $\mutrue$ with probability at least $1-\delta$
  (\Cref{thm:hw-decomp}).  The proof reduces to two independently
  checkable invariants---mass-conservative refinement
  (\Cref{lem:mass-conservation}) and SMT-certified closure
  (\Cref{lem:closure-soundness})---plus an application of an
  anytime-valid concentration inequality per leaf.

\item \emph{Anytime-valid certification via PrPl-EB.}  We give the
  first application of the predictable-plug-in empirical-Bernstein
  confidence sequence in the context of program reliability
  estimation.  The bound is closed-form, variance-adaptive, and
  requires no prior knowledge of per-leaf variance, strictly
  tightening standard Hoeffding and Wilson constructions under
  adaptive stopping.

\item \emph{Implementation and evaluation.}  We implement \toolname{}
  as \toolname{}, a Python prototype with a Z3 backend and a benchmark
  suite of $12$ integer programs.  On rare-event benchmarks whose path
  structure aligns with the input distribution, \toolname{} reduces
  sample counts by up to two orders of magnitude relative to plain and
  stratified Monte Carlo.
\end{enumerate}

\paragraph{Paper organisation.}
\Cref{sec:problem} formally states the problem and defines the input
distribution class.  \Cref{sec:related} surveys related work.
\Cref{sec:algorithm} presents \toolname{} and proves
\Cref{thm:hw-decomp}.  \Cref{sec:implementation} describes the
\toolname{} implementation.  \Cref{sec:results} reports experimental
results.  \Cref{sec:conclusion} concludes.
